\chapter{Validating the measurement}
\label{app:validation}
\emergencystretch=3em

This appendix records how the measuring instruments were checked and corrected. Nothing in it is a result of the analysis. It shows that the figures of Chapter~\ref{ch:results} come from instruments whose own failures were looked for and found. Section~\ref{app:repairs} lists every correction in one table, and the sections after it give the detail behind those that need more than a table row.

\section{History of the corrections}
\label{app:repairs}

uigraph was developed within this thesis, and three of its versions are named in this appendix, each by one label.
\begin{itemize}
\item The \emph{early development version}: commit \EarlierCommit{}, before the corrections the evaluation
  prompted. On it k8s dashboard had nodes and no edges (Table~\ref{tab:defects}).
\item The \emph{intermediate version}: extractor \RoundExtractor{}, with its
\RoundMustEdges{}-edge must tier and the tables as printed at \RoundCommit{}.
\item The \emph{final version}: extractor commit \texttt{c42f220}, harness 1.9.0 and mutation
harness 1.9.1. The demotion-rule sweep and the may-reason census were recorded at
\texttt{d05227c}, which adds to it only a measurement switch that leaves the graphs unchanged.
\end{itemize}
The chapters present only the final version, whose tables of Section~\ref{sec:res-round2} were regenerated at its commit.
Of the intermediate version only the must-tier figures are cited (Table~\ref{tab:repairs}), as printed at \RoundCommit{}; the runs behind
the printed tables were lost, and an archived re-run of the same extractor gave a different loss ledger.
The figures of the early development version and of the intermediate must tier are historical, measured with
\texttt{eb9dab80} and before and read from git at \texttt{f17a302} and \texttt{7d109c7}.

Table~\ref{tab:repairs} records how the instrument got to the corrected one of Chapter~\ref{ch:results}: what was wrong, how it was found, what was changed, and what the change did to the reported figures. The figures of superseded instruments cannot be regenerated with the harness this thesis ships. Each is read from a named commit's own generated results or from an earlier run that was not kept, and in the second case the row describes the effect in words. The reviews and red-team panels named in the table were made with the AI assistant that the declaration at the front of this thesis describes, and every defect they reported was reproduced against the code before it was repaired.

\begingroup\footnotesize\setlength{\tabcolsep}{2.5pt}
\begin{xltabular}{\textwidth}{@{}rYYYY@{}}
\caption{The corrections made to the instruments, in the order the chapters record them.
\emph{Found by} gives the method as the main text names it.}
\label{tab:repairs}\\
\toprule
No. & What was wrong & Found by & What changed & Effect on the figures \\
\midrule
\endfirsthead
\toprule
No. & What was wrong & Found by & What changed & Effect on the figures \\
\midrule
\endhead
\bottomrule
\endfoot
\multicolumn{5}{@{}l}{\emph{Protocol and harness (Section~\ref{sec:corrections})}} \\
1 & Test-source cross-check: a URL parameter named in a test could match a literal route segment, which gave matches that were not real screens.
  & Seen as the recall figure of one application falling when the rule was tightened.
  & A parameter no longer matches a literal segment. A negative test, which names a URL that must not exist, is excluded.
  & The earlier recall figure cannot be regenerated. The end-to-end recall of that application is now \EtoERateAlfresco{}. \\
2 & The first witness check reported a low precision for one application.
  & Inspection of the flagged edges. All were may edges of confidence one half, produced by one over-approximated call site connected to every declared route.
  & No extractor change is recorded. Over-approximated sites are reported apart, since the check claims no target agreement for them.
  & The low figure described the measurement, not the extractor. The must tier of the same run had no failing edge. \\
3 & Mutation testing: a mutant whose target is absent from the baseline graph repeats a recovery gap already counted elsewhere.
  & Not stated beyond the misses of one application.
  & Such mutants moved to a separate category.
  & Most of that application's misses left the detection rate. \\
4 & The first construct denominator counted the substring \texttt{routerLink}, which also matches \texttt{routerLinkActive} and its options.
  & A construct count that tripled, a figure that looked implausible.
  & Constructs are counted by the independent route parser that also produces the loss breakdown, so ratio and breakdown share one definition.
  & For alfresco the independent route parser finds \ConstructsAlfresco{}. The substring count was about three times larger. \\
5 & The first independent route parser did not look inside wrapper calls, did not follow a lazily loaded module to the routing module it imports, and walked lazy tables without their prefix.
  & Coverage that reached 100 per cent: every declared screen of mifos had a node.
  & The independent route parser reads wrapper calls, follows the imported routing module, and keeps the prefix.
  & mifos has \RoundTwoEarlierOracleMifosDeclared{} declared screens, of which \NodeAuditMatchedMifos{} have a node. The extractor was the same in both runs. \\
6 & Must conversion was first defined as must edges divided by constructs.
  & A ratio that rose implausibly after attribution (mifos).
  & The ratio is computed per construct, and edges whose screen set is a whole subtree are stored once (Section~\ref{sec:attribution}).
  & The rise came from how edges were stored, one per screen, and not from evidence. The ratio is defined in Section~\ref{sec:metrics}. \\
7 & The mutation harness patched the first textual match of a construct, judged detection on node and edge sets alone, and retargeted only the head of a commands array.
  & A review of the harness, and a detection rate that moved by more than twenty points.
  & Patches are made at the construct's offset. Detection is judged on the witness set as well. A retargeting mutant points at another declared route, a guard mutant on a screen no edge enters is excluded, and a rename counts when the renamed route appears.
  & The detection rate for realworld rose sharply. The corrected rule was applied with the same seed but not the same sample: \MutEarlierSampled{} mutants in the earlier run, \MutSampled{} in the final one, \MutOverlap{} at the same location (operator, file and line; destinations may differ). Of the \MutEarlierFailures{} earlier misses, \MutEarlierFailuresAbsent{} are not in the final sample. The earlier run showed six paired disagreements in favour of attribution, five of them retarget mutants whose operator replaced only the head of the commands array; with the operator corrected the baseline detects those five too. \\
8 & The guard exclusion rule found a route by the nearest path literal before \texttt{canActivate}, which fails for routes with children, empty paths, constants and template literals.
  & Reading the graph at the sites the rule had let through.
  & The route is found by matching the braces of the enclosing object over source with comments and strings blanked. A route yielding no screen is excluded.
  & Four guard mutants that cannot be observed had been scored as misses, and the discussion had explained them as a defect in how the graph represents a guard. They are now excluded. The superseded run sampled \MutEarlierSampled{} mutants. \\
9 & Removed guards were judged wrongly in two ways: an edge between two children of a guarded parent no longer carried the parent's guard, and a comparison of sets read one of three identical guard texts as unchanged.
  & A change in the extractor (it gated an edge only on the guards its navigation runs), not a figure.
  & A removal is detected when some edge loses a guard, and excluded when no baseline edge carries it. Guards are counted, not collected.
  & Which mutants are sampled is unchanged. How a sampled mutant is scored changes. \\
10 & The construct denominator took a construct from any line its pattern matched, including a commented-out link and a sentence of a code comment.
  & A table script stopped: the loss classifier found no reason for a commented-out link in mifos.
  & Comments are blanked before parsing, with the compiler's own scanner for TypeScript.
  & One construct leaves the denominator in each of bitwarden and mifos, and nothing else changes. \\
\addlinespace
\multicolumn{5}{@{}l}{\emph{The extractor and the deterministic tier (Section~\ref{sec:res-round2})}} \\
11 & A guard on a route that a wrapper function returns was not read, so everything below it was published as unconditional.
  & Red-team panels reading the graph and the source after commit \EarlierCommit{}. The witness check did not find it.
  & The guard set of a screen also reads routes a function returns.
  & At \EarlierCommit{} the tier held \EarlierMustEdges{} must edges with \EarlierMustViolations{} witness failures, read from that commit's own results. The failures were target mismatches on edges the unread guard left on the deterministic tier. Reading the guard cut mifos's tier to \RoundMustEdgesMifos{} on the intermediate version. \\
12 & A \texttt{canActivate} element that is a call, such as a guard factory, was skipped, so a route guarded only by one read as unguarded.
  & Red-team panels, as in row 11.
  & An element the reader cannot name gates the edge anyway.
  & bitwarden's deterministic edges into its root no longer read as unconditional. \\
13 & The witness check opened only the first of several constructs that yield one edge.
  & Correcting the check itself, so that it judged every cited site. This found edges in bitwarden whose later citations targeted \path{/vault}, which the extractor had not read.
  & The check opens every cited site. An edge took the weakest modality over all constructs citing it, a rule replaced in row 18.
  & Later citations are \WitnessMergedCitationShareAll{} of cited sites, and \WitnessMergedCitationShareMax{} in the largest application. The intermediate tier then held \RoundMustEdges{} must edges, with \RoundMustClean{} passing and \RoundMustViolations{} failures. \\
14 & The demotion rules were incomplete. A link under \texttt{*ngIf}, inside \texttt{@if}, in a callback, in a \texttt{catch} clause, or after \texttt{await} or an early return was emitted as must.
  & A review of every construct cited by the intermediate tier's \MustCtxEdges{} edges (\MustCtxDistinct{} distinct constructs, cited \MustCtxConstructs{} times), then red-team panels. The witness check cannot find this defect: it compares destinations and does not ask whether a construct runs every time.
  & A site is demoted under every condition of Section~\ref{sec:model}. The panels extended the rules to members no handler calls unconditionally, components mounted under a template condition, ICU cases and addresses two routes can answer; a later review extended them to code after a navigation that may start another.
  & \MustCtxConditional{} of the \MustCtxConstructs{} citations ran under such a condition. Table~\ref{tab:mustledger} records the condition read at each construct. \\
15 & A guard on an empty-path child without a component, which Angular matches on every navigation to the parent, was not read. \texttt{canMatch}, \texttt{canActivateChild} and \texttt{canLoad} were not read at all.
  & The review of row 14.
  & Such a guard gates the screen above it, and the other three guard kinds are read.
  & \MustFateGuard{} of the intermediate tier's edges, all of bitwarden's, are guarded. \\
16 & The repair of row 15 gated an edge on every guard of its destination, which demotes navigations the router lets through untouched.
  & A red-team panel.
  & The extractor computes, for each source and destination, the guards the router runs: levels left, matched, entered and kept.
  & Navigations between screens that share a guarded parent are let through, and mifos holds \CensusMustMifos{} must edges again. \\
17 & A navigation to the address its source already displays was treated as one that occurs on every execution.
  & The review of row 14.
  & A shell or anchor edge whose destination is one of its screens is drawn per screen: must from the others and a may self-loop at the destination.
  & \MustFateSelfLoop{} of the intermediate tier's edges is such a self-loop. \\
18 & The weakest-modality rule of row 13 discarded a certainty the edge had. realworld's Global Feed link runs whenever its page does, yet two other links to the same root do not.
  & Reading the edges that survived.
  & An edge takes the modality of its most certain construct, which becomes its witness. Every other construct records the modality it earns itself, and the witness check judges the must tier only on constructs that earn must.
  & The final tier is \CensusMust{} edges resting on \CensusMustSites{} constructs. It is a different set from the intermediate tier, not that tier made smaller. Superseded tables counted every construct a must edge cites as converted; the current ones count a construct only where it earns must itself. \\
\addlinespace
\multicolumn{5}{@{}l}{\emph{The independent route parser and the denominators (Sections~\ref{app:readers} and~\ref{app:setmoved})}} \\
19 & The independent route parser did not read a path written through a constant, an enum or \texttt{\$localize}, a first-party import through \texttt{baseUrl}, or a route array named rather than written.
  & Investigating the disagreements of the first node audit (Section~\ref{app:readers}). Each gap had scored the extractor against paths the independent route parser itself had not read.
  & The independent route parser reads all three forms.
  & Before the last of them, the independent route parser saw almost none of kubernetes/dashboard's declared screens and scored nearly every node the extractor minted there as unmatched. \\
20 & The independent route parser read Storybook fixtures as declared screens, read a sibling application's tables reached through a tsconfig alias as this application's, and recorded a path in an \texttt{Object.freeze} table as unresolvable.
  & The dead-link corrections (Section~\ref{app:deadlinkfixes}).
  & The independent route parser sets Storybook files aside, ignores the sibling's tables, and reads the frozen table. Dropping the frozen routes from the denominator as unreadable was considered and rejected, since they were readable one call away.
  & The declared sets of \RoundTwoSetMoved{} of the \CensusApps{} applications moved, and in \RoundTwoSetMovedCountSame{} of them the count stayed the same (Section~\ref{app:setmoved}). \\
21 & A file produced by an interim revision of the independent route parser, dropped into the directory the tables are generated from, was read as the earlier independent route parser's output. The script that produces the tables could not detect this, and an absent file was read as a denominator that did not move.
  & Fixing the provenance of the comparison.
  & The earlier independent route parser's own committed output is read from the repository.
  & The bitwarden decomposition built on that file is withdrawn; Section~\ref{app:setmoved} gives the corrected one. \\
22 & The first-audit extractor reached route tables without their parent prefix, and minted a node at \texttt{*} for each catch-all.
  & Resolving the disagreement of the first node audit (Section~\ref{app:readers}).
  & The repairs of the loss list: a table not reached, a path named and not written, a destination named and not written (Section~\ref{sec:res-defects}).
  & Of \RoundTwoFirstRoundExtra{} nodes with no declared route, \RoundTwoFirstRoundTruncated{} were the tail of a declared route and \RoundTwoFirstRoundSynthetic{} were catch-alls. On the final version \RoundTwoExtraNodes{} nodes match no declared screen, \RoundTwoExtraNodesWildcard{} of them at \texttt{*}. \\
23 & The grounds that exclude a construct from the denominator were tested after the graph had been consulted.
  & Not stated.
  & The grounds are decided from the construct's own text before the graph is consulted. The comparison refuses to print if a construct changes its ground between the two runs.
  & The after column had been a ratio over more sites than the before column. Both are now over the same \RoundTwoInScope{} sites. \\
\addlinespace
\multicolumn{5}{@{}l}{\emph{Findings of the runtime route oracle (Section~\ref{sec:res-probe})}} \\
24 & A route that declares a named outlet could take the primary address from a sibling. The routes beneath a named-outlet route were also minted at the primary address.
  & The runtime route oracle returned two alfresco addresses that the framework answers with its catch-all. Both static readers had read the subtree alike, so coverage of 100 per cent could not expose them.
  & Outlet context is carried down the walk, and no screen is minted for a subtree beneath a named outlet.
  & On the compared graph the extractor declares \ProbeInvented{} screens at addresses the framework mounts nothing at. What remains is a limit on recall (Section~\ref{sec:named-outlet}). \\
25 & A lazy route that names a module through a barrel, which re-exports a feature module that imports the routing module, was followed one hop only. The table was then scanned at the root.
  & The runtime route oracle, compared with the route set as the static readers held it before. Both readers had arrived at the defect independently.
  & Both follow the chain to the table, bounded, and refuse at each hop a file the route sits in, a file another route loads, and a file whose tables are registered at the root.
  & \ProbeBeforeNotMounted{} addresses in \ProbeBeforeWhere{} were wrong, the two alfresco ones of row 24 among them. Node coverage had been 100 per cent against a denominator carrying the same wrong addresses. \\
26 & The test that keeps a lazy route from adopting the application's root table read \texttt{RouterModule.\allowbreak forRoot}, so it fired nowhere in the corpus.
  & Instrumenting the predicate, at commit \texttt{d40c8dc}.
  & Both tools mark root tables per array, taken from the argument of either registration API.
  & A guard that never fires had been credited with a property nothing showed it to have. The counts are in Section~\ref{sec:ownership}. \\
\addlinespace
\multicolumn{5}{@{}l}{\emph{The dead-link census (Section~\ref{app:deadlinkfixes})}} \\
27 & An earlier correction followed only the redirecting kind of catch-all route.
  & Attempts to refute the dead-link finding.
  & The extractor follows both kinds, holds the edge at may, and records where the catch-all sent the URL.
  & The correction had held for the application it was written against and failed for the two that needed the other kind. \\
28 & Resolving a target and calling a link dead were conflated. bitwarden's authentication paths went through an unread enum, so the wildcard forwarded every one of them.
  & Attempts to refute the finding.
  & The enum and the wrapper call are read. Each application records a note when its route set is a lower bound. Each caveated candidate is put to the framework's own route configuration.
  & The wildcard had raised that application's recovered figure. Candidates withdrawn: \RoundTwoFrameworkMountsIt{} because the framework serves a screen, \RoundTwoFrameworkPrefix{} because the matching route places none, \RoundTwoFrameworkRedirects{} because a redirect answers. \RoundTwoWildcardCaveated{} keep the caveat. \\
29 & Routes the tools could not read were taken as routes the application never declares: a route array built by a spread, ghostfolio's \texttt{matcher} property, and alfresco's tsconfig \texttt{paths} alias.
  & Attempts to refute the finding.
  & Both tools record the holes they cannot read, and no claim that a path is undeclared is counted inside one.
  & ghostfolio's \path{/admin/users} had been reported as a confirmed dead link, and two live links of alfresco as dead. All three constructs are read now. \\
30 & The route-set test read only the first quoted segment of a commands array and so tested a prefix as if it were the target.
  & Attempts to refute the finding.
  & The whole array is read as a pattern, computed segments as holes, and matched against the declared routes pattern to pattern.
  & It cleared alfresco's two candidates and found further sites in mifos. The census now holds \RoundTwoDeadLinkCandidates{} candidate sites. \\
\addlinespace
\multicolumn{5}{@{}l}{\emph{Mutation sweep and end-to-end check}} \\
31 & On its first run in the final revision, the sweep of the demotion rules found rules the router fixes had added without a test.
  & The mutation sweep, at commit \SweepEarlierCommit{}.
  & Tests were written for most survivors. The others were retired as equivalent, since no input can tell them from their absence.
  & \SweepEarlierSurvived{} of \SweepEarlierMutants{} mutants survived; \SweepEarlierCovered{} now have tests and \SweepEarlierRetired{} were retired. At \SweepCommit{} every one of \SweepMutants{} deletions fails some test. Writing the tests exposed a relative link that climbs above the root. \\
32 & The end-to-end scorer compared a redirected navigation on the landed address, against the preregistration.
  & A review panel (4 Nov), checking the scorer against the preregistration; fixed at commit \texttt{0172867}.
  & The scorer applies the preregistered rule for redirected navigations.
  & mifos now scores \EtoEMifosHit{} of \EtoEMifosPairs{} pairs; \EtoEMifosRedirect{} pairs were redirected, and \EtoEMifosHitRedirect{} of them hit. The earlier figure, scored on the landed target, was \EtoEMifosHitAfterRedirects{} of \EtoEMifosPairsAfterRedirects{}. \\
\end{xltabular}
\endgroup

Five of the protocol corrections (rows 2, 4, 5, 6 and 7) were found by a figure that looked implausibly good or bad. Each turned out to be a property of the measurement and not of the extractor. A surprising figure is therefore checked against its inputs before it is written down.

\section{Nodes the two static readers disagreed about}
\label{app:readers}

This section records the direction of the differential comparison in which the extractor emits nodes for which the independent route parser holds no declared screen, as the first node audit measured it, with both tools as they stood before the repairs of Table~\ref{tab:repairs}. That audit is historical: harness 1.2.0, recorded extractor \texttt{4505fd2a}, absent from the repository, committed table only. Section~\ref{sec:res-round2} repeats the comparison with both tools corrected, and the figures change considerably.

In the first audit the independent route parser held no declared screen for \NodeAuditUnmatchedAll{} of the \NodeAuditGraphNodesAll{} nodes in the corpus. For \NodeAuditBlindAll{} of them it could not follow the route that would have declared them, which is its own limit (\NodeAuditBlindKdash{} in k8s dashboard). For \NodeAuditOutsideSAll{} it found a route of a kind the model excludes from the screen set, a path prefix or a redirect. The remaining \RoundTwoFirstRoundExtra{} had no declared route of any kind at the node's path, concentrated in mifos (\NodeAuditNoRouteMifos{}) and bitwarden (\NodeAuditNoRouteBitwarden{}). This third group was the largest unexplained behaviour of the first audit, because the two methods could not adjudicate it.

The later repairs resolved it. \RoundTwoFirstRoundTruncated{} of the \RoundTwoFirstRoundExtra{} sat at the tail of a route the independent route parser had declared in full, because the extractor had reached the table without its parent prefix. The other \RoundTwoFirstRoundSynthetic{} were catch-all routes at the path \texttt{*}, which the independent route parser's policy keeps out of the screen set, so there the two tools disagree by convention and not over a table either missed. The repairs of Table~\ref{tab:repairs} were accordingly almost all in the extractor.

On the instrument of Chapter~\ref{ch:results}, \RoundTwoExtraNodes{} nodes remain that no declared screen matches, and \RoundTwoExtraNodesWildcard{} of them sit at \texttt{*}. The independent route parser declares a screen for every other node the extractor mints, and the runtime route oracle confirms that the framework mounts a component at every address the extractor claims. This section therefore records a past state and not a standing defect, and the reason a third, independent reader was added.

\section{How the declared-screen set moved}
\label{app:setmoved}

Table~\ref{tab:coverage} gives node coverage, and \RoundTwoFullNodes{} of the \CensusApps{} applications now hold a node for every declared screen. The denominator of those ratios needs separate treatment, because the declared-screen \emph{count} is a poor summary of the declared-screen \emph{set}. Between the independent route parser before the repairs and the one reported here, the declared sets of \RoundTwoSetMoved{} applications changed, and in \RoundTwoSetMovedCountSame{} of them the count stood still. In ghostfolio the count did not move (\RoundTwoEarlierOracleGhostfolioDeclared{} before and \RoundTwoEarlierOracleGhostfolioNow{} after) while every address changed, since the earlier parser recorded a route's path as a verbatim expression. The largest move is k8s dashboard, from \RoundTwoEarlierOracleKdashDeclared{} declared screens to \RoundTwoEarlierOracleKdashNow{}: the earlier parser could not resolve its \NodeAuditBlindSpotsKdash{} lazy targets and declared only the screens of two eagerly registered tables, so coverage against that denominator was 100\,\% of almost nothing, a fact about the instrument and not about the application. In bitwarden, \RoundTwoEarlierOracleBitwardenStorybook{} Storybook fixtures left the set, and the moved screens are the same screens at a different address, which a diff of the set shows and a count does not. A moving denominator can distort a coverage figure unseen, and one that keeps its size can do so too, so every part of the subtraction is generated and can be checked. It is made against the earlier independent route parser's own committed output, read from the repository, because an interim revision's file was once read in its place (Table~\ref{tab:repairs}, row 21).

\paragraph{Witnesses outside the application directory.} The independent route parser and the extractor agree about an edge only where both read the same source. \RoundTwoOutsideWitness{} edges across \RoundTwoOutsideWitnessApps{} applications (\RoundTwoOutsideWitnessWorstN{} of them in \RoundTwoOutsideWitnessWorst{}, \RoundTwoOutsideWitnessWorstMust{} of those on the deterministic tier) have their only witness in a file outside the application directory, such as a shared library or, in bitwarden, the repository's other front ends. Both tools read those files, so their agreeing there is not independent evidence. The figure is reported and the edges are not excluded, because a precision figure is a claim about the edges the tool emitted, wherever it read them, and a reader who wants the figure over the measured directory alone can subtract the count. Whether a sibling front end's components are reached in the application under measurement is not established here. The dead-link census of Section~\ref{sec:deadlinks} decides the other way, because calling a link dead asserts something about the application, so there the boundary decides.

\paragraph{A coverage ratio cannot see a screen nobody can address.} Some of bitwarden's declared screens were for a time outside the denominator, because their paths are assembled from a constant that neither tool resolved, and coverage was then 100\,\% with no sign of the error. The only defence is to report what was set aside, which is why the blind-spots column is printed beside every coverage figure in Chapter~\ref{ch:results}. That \RoundTwoFullNodes{} applications reach every declared screen is accordingly a weaker result than it sounds: it shows that the walk finds the route tables, not that the graph is complete.

\section{Every remaining site names a reason}
\label{app:reasons}

\RoundTwoUnexplained{} sites are neither recovered nor accounted as dynamic or unplaced. A script records a reason for each and exits non-zero if any site reaches its end without one, so a site the analysis cannot explain is a failing state and not a silent remainder. The script enforces that a label is recorded, not that it is correct. The check covers these sites only: the remaining attribution losses of Table~\ref{tab:remaining} include \AttrRemainingUnexplained{} recorded as having a scope and a destination with no reason found, and that label is accepted rather than failed.

The largest group, \RoundTwoUnplacedClass{}, is a class the extractor parsed and found a navigation site in but could place under no screen, because no relation the implementation builds reaches it (Sections~\ref{sec:inputs} and~\ref{sec:approximations}). The reason given is the extractor's own note, a claim about the tool and not about the application. A second group, \RoundTwoTargetNoRoute{} sites whose target matched no route, is split by testing the target against the independent route parser's route set: \RoundTwoRelativeUnresolved{} are relative targets with no base, \RoundTwoUnderUnreadable{} lie in a subtree one tool could not read, \RoundTwoNoRouteCaveated{} are in an application whose route set is a lower bound, and \RoundTwoDeadLinkUnanswered{} are accounted for by nothing but the application.

\section{The four dead-link corrections}
\label{app:deadlinkfixes}

Four corrections belong with the dead-link finding, and all four came from attempts to refute it (Table~\ref{tab:repairs}, rows 27 to 30). Three are the same mistake made in three places: from a navigation site, a route neither tool can read looks like a route the application never declares.

\paragraph{Catch-all routes.} A target that no route declares does not necessarily go nowhere: a \texttt{**} route catches every unmatched URL, either by forwarding it with \texttt{redirectTo} or by rendering a \texttt{component}. The extractor follows both, holds the edge at may, and records where the catch-all sent the URL. An earlier correction followed only the redirecting kind, and failed for the two applications that needed the other kind.

\paragraph{Resolving a target is not calling it dead.} bitwarden declares its authentication paths through an enum the extractor had failed to read, so the wildcard forwarded every one of them and raised the application's recovered figure. The enum is read now, and the extractor records one note per application saying when its route set is a lower bound. Each caveated candidate is also put to the route configuration the framework holds (Section~\ref{sec:res-probe}). \RoundTwoFrameworkMountsIt{} are withdrawn because the framework serves a screen at the target, \RoundTwoFrameworkPrefix{} because the target is matched by a route that places no screen of its own, and \RoundTwoFrameworkRedirects{} because a redirect answers it. \RoundTwoWildcardCaveated{} sites keep the caveat.

\paragraph{Holes counted as absence.} A route array built by a spread was not followed, ghostfolio mounts its user administration screen with Angular's \texttt{matcher} property that neither tool read, and alfresco loads children through a tsconfig \texttt{paths} alias that neither tool resolved. An earlier state of the census reported ghostfolio's \path{/admin/users} as a dead link confirmed by hand, but the hand check had searched for \texttt{path: 'users'}, which is not how that route is spelled, and it reported two live links of alfresco as dead. All three are read now. Both tools record the holes they cannot read, and no claim that a path is undeclared is counted inside one.

\paragraph{A prefix tested as the target.} The route-set test read only the first quoted segment of a commands array, so \texttt{['/nodes', id, 'rules']} was tested as \texttt{/nodes}. The whole array is now read as a pattern and matched against the declared routes pattern to pattern.

\paragraph{What the candidates are.} The dead-link candidates are \RoundTwoDeadLinkCandidates{} sites at \RoundTwoDeadLinkLocations{} places in the source, \RoundTwoDeadLinkUrls{} distinct URLs, in \RoundTwoDeadLinkApps{} applications (\RoundTwoDeadLinkWhere{}). The count is of sites and not of defects: one wrong prefix written into a shared template is one mistake and several sites. Every other category of the \RoundTwoWildcardRescued{} sites that a catch-all route answers is a limit of one of the two tools and not a defect of an application, so the residual is what is left once the tools no longer blame the applications for their own blind spots.

The result is a candidate list and not a verdict. The corrections show that the tools' blind spots are not yet all known, and a route that both tools fail to read looks, from outside, like a route that does not exist. Each candidate was checked by hand against the application's own route files, which is how the \texttt{appusers}, \texttt{manage-jobs} and \texttt{configurations} mismatches of Section~\ref{sec:deadlinks} are quoted, and its target was put to the route configuration the framework holds, read from the running application (Section~\ref{sec:res-probe}). The classification used for that premise was itself held against what the framework does, on \ProbeAnswerChecked{} navigated addresses, and agrees on \ProbeAnswerAgreed{}. Still unchecked is the other premise, that the link is reached at run time at all: no candidate was checked by navigating to it in a running application, and a target is matched against the framework's patterns and not navigated to one address at a time. What the method contributes is the list, with a file and a line each, instead of leaving the candidates to be met at run time.

\section{Asking the running framework}
\label{sec:res-probe}

Every result so far rests on two readings of the same route files, and where both tools read a route table at the wrong prefix they agree on a wrong address and the witness check passes (Section~\ref{sec:protocol}). This section reports a third reading, which does not read the text at all.

\subsection{How the framework was asked}

Each application was built in its development configuration and loaded in a headless browser. Once it had bootstrapped, the live \texttt{Router} was found through Angular's debug surface, its route configuration read directly, and every lazily loaded subtree expanded by calling the application's own \texttt{loadChildren}. Nothing in the procedure parses a route file, so a disagreement with both static readers cannot be an artefact of a reading they share. No server was run, since the method needs the application to boot, not to work. An address is counted as mounted when the router's configuration places a component there on the primary outlet, whether or not a guard would admit a visitor. All \ProbeApps{} applications bootstrapped and \ProbeRoutes{} route entries were read. \ProbeLazyFailures{} lazily loaded subtrees were left unexpanded, \ProbePageErrorApps{} applications raised an uncaught error, and \ProbeConsoleErrorApps{} logged errors to the console, mostly failed requests to the server that is not there, which does not bear on the route configuration. \ProbeMatcherRoutes{} route in the corpus decides its address through a \texttt{UrlMatcher}; it and anything beneath it have no address any of the three readers can resolve to a string, so they are outside both sides of the comparison. Table~\ref{tab:probe} reports the result.

Which entries answer an address with a screen is itself a reading, and it is checked against what the framework does. The harness navigates a running application to an address and records the recogniser's verdict, the event Angular emits when a URL has matched. Over \ProbeAnswerApps{} applications and \ProbeAnswerAddresses{} addresses, chosen to include every address this chapter names, the classification taken from the configuration agrees with the navigation on \ProbeAnswerAgreed{} of the \ProbeAnswerChecked{} it decides, and declines to decide \ProbeAnswerUndecided{}, the address the \texttt{UrlMatcher} route claims. The comparison of Table~\ref{tab:probe} is held against the same navigations on the two implications that bind it, and agrees on \ProbeCmpAgreed{} of \ProbeCmpChecked{}. The addresses are those the text reasons about and not a sample of the corpus, and the \ProbeMissed{} screens the extractor does not mint are not individually checked, and are not claimed to be. The comparison turns on two answers: an address whose route carries children places no screen of its own (mifos's \path{/collections} answers with the application's chrome and an empty outlet), and in an application with a catch-all the configuration mentioning an address is no evidence that anything is mounted there.

\subsection{What the framework mounts and the extractor does not}

The two directions are not symmetric. A screen the extractor does not mint costs recall, since the graph says less than it could. An address it mints that nothing mounts would be a false claim about an application. The corpus holds \ProbeMissed{} of the first and \ProbeInvented{} of the second.

Of the \ProbeMissed{}, \ProbeMissedGhostfolio{} are one construct in ghostfolio: product comparison pages generated from a list of competitor records spread into the table, which the extractor declares a route source it could not read. This shows that the cost of a declared blind spot can be far larger than a count of notes suggests. \ProbeMissedMifos{} are one construct in mifos, a choice between two route tables by a constant computed from \texttt{window.env} at deployment time, which no reader of the code can resolve, so the extractor declares the subtree unread. \ProbeMissedAlfresco{} are in alfresco and no two share a cause: two are outside the directory the measurement scopes to, and \path{/about} is declared in a JSON extension manifest that no reader of the application's code can see, and that the runtime route oracle sees because the framework assembles the table from the manifest. The remaining \ProbeMissedBitwarden{} is bitwarden's vault, whose route is produced by a helper that chooses between two components by feature flag. Each loss is declared, which is the trade-off this analysis is designed to make. The mifos and alfresco constructs are recorded in the approximations of Section~\ref{sec:approximations}, and alfresco's route set is a declared lower bound on other grounds.

\subsection{Nothing is missing that was not declared}

\ProbeMissedWithoutDeclaration{} applications hold a screen the framework mounts and the extractor has no node for without the extractor having said that part of its route set was beyond it. Every application in Table~\ref{tab:probe} that has a missing screen is an application that declared one. The correspondence is per application and not per cause, which is the weaker reading and the only one measured here: alfresco declares that it could not read ten spreads, yet its missing \path{/about} has nothing to do with a spread. The measurement supports the claim that no application is missing a screen while asserting that its route set is whole, and not the claim that every missing screen is behind the construct its application named. On this corpus the correspondence also holds the other way: \ProbeDeclaredButComplete{} applications declare part of the route set unread and are missing nothing, and \ProbeSilentAndComplete{} declare nothing and are missing nothing. That is not a property of the analysis and is not offered as one.

This changes what the dead-link census of Section~\ref{sec:deadlinks} rests on. The census used to suppress every candidate from an application whose route set is a declared lower bound, with a per-application licence lifting the suppression where this comparison found the declared set complete. On the figures above the licence lifts it for \RoundTwoProbeLifted{} applications, since both applications the census draws from are missing a screen of their own (mifos \ProbeMissedMifos{} and bitwarden \ProbeMissedBitwarden{}). Section~\ref{sec:deadlinks} therefore asks the question per candidate, of the same framework configuration.

When this check was first run, the correspondence failed: alfresco and bitwarden, whose route tables are built by a helper call spread into the table, asserted complete route sets while holding unread route sources, and bitwarden's silence let its sites enter the dead-link census on a route set that was missing a screen. Both readers now record a spread they cannot follow as a route source they did not read.

\section{What a route table's evidence does not establish}
\label{sec:ownership}

One assumption underlies every approximation in the reader of route tables: a route table that no site in the application claims (no \texttt{loadChildren} names it, no spread or \texttt{children} carries it, no registration call takes it) is read as a table the application mounts at its root. It is an over-approximation, since a table registered nowhere the compiled build reaches is mounted at no address at all, and reading it at the root puts its entries into the graph as screens. It is taken deliberately, because the alternative loses the routes of every application whose root registration neither tool recognises. The same holds when the only site that loads a table lies in a test specification, which occurs nowhere in this corpus.

Three further approximations in that reader share one shape, and each was found by attacking a repair of the previous one. All three are declared in the tool's own notes, and none was closed, because closing it would refuse a construct the corpus actually uses. First, a \texttt{Routes}-typed array shows only that its author declared it to some router, so a sibling application's typed table reached through a workspace alias is still elected as the owner of a lazy route array, and still minted under the lazy prefix with a deterministic edge into it; requiring the owner to lie inside the application directory would also refuse a route array that legitimately lives in a shared library, a pattern the corpus does use. No corpus application distinguishes the two, so the gap is left open and declared. Second, the \texttt{loadChildren} specifier is not bounded, so a loader whose import resolves into another project's alias root mounts that project's table under this prefix; a guard against it was applied experimentally, turned several declared screens of one application into declared blind spots and lowered its coverage, and by path alone the analysis cannot tell a shared lazy library from another application's route table. Third, a routing module imported by two lazily loaded modules and loaded by no route is read as a table of the project root, which costs the affected routes their prefix.

\paragraph{The test for a root table.} A fourth case was found by asking what the test that keeps a lazy route from adopting the application's root table actually read. The test looked for the candidate file's own call to \texttt{RouterModule.forRoot}, which made it sound where it fired, and is also why it fired nowhere: an application that registers its root with \texttt{provideRouter} is not recognised at all, and one that declares its root array in one file and registers it in another is recognised only in the registering file. Instrumenting the predicate turned the limit into a measurement. When this repair was made (commit \texttt{d40c8dc}, before the final extractor, and not re-measured on it), the test was asked 698 times across the six applications in the independent route parser and 2798 times in the extractor, and answered yes on none of them. These counts were read from an instrumented run whose output was not committed, so they cannot be regenerated. The independent route parser's own source comment gives 2798 for its predicate, so which count belongs to which tool is not settled. A guard that never fires is worse than no guard, because a commit message and a chapter credit it with a property nothing has shown it to have. Both tools now mark root tables per \emph{array}, taken from the argument of either registration API.

\paragraph{Repairs that stand on fixtures only.} Five repairs of the same reader stand on constructed fixtures and on nothing else. Each produced an address the application does not mount, or a declared loss where nothing was lost, and each is now refused. None occurs in this corpus, so every figure in Chapter~\ref{ch:results} is identical with and without them, and no result here is evidence that they work; their evidence is the fixtures, which fail when the rule each pins is removed. The same holds of the refusals in the independent route parser, such as the test that stops a lazy route adopting a table another route owns, which answers yes in neither tool on the corpus. A rule the corpus exercises and a rule only a fixture exercises are different kinds of claim, and a corpus of six applications is not a sample that makes an analysis sound: five constructs it happens not to contain were enough to break two of the properties this thesis claims.

\paragraph{The named outlet.}
\label{sec:named-outlet}
A route carrying \texttt{outlet: 'name'} contributes no screen on the primary outlet. The extractor refused such a route from the start, but where a named outlet opened a \texttt{children} subtree the routes beneath it were still minted at the primary path. alfresco mounts its document viewer that way, and the graph held \path{/view/:nodeId} and \path{/search-libraries/view/:nodeId}, two addresses that application serves only through the viewer outlet. The independent route parser read the subtree the same way, so coverage was 100\,\% over a set holding both, and no comparison between the two tools could show the gap. The case was first written down as open, on the reasoning that an approximation stated is one a consumer can price. That reasoning was wrong, as the runtime route oracle of Section~\ref{sec:res-probe} showed: a navigation to that address in the built application is answered by the catch-all route, and the must-tier property admits no exception on the strength of a written caveat. The extractor now carries outlet context down the walk and mints no screen for such a subtree. What remains is a limit on what the graph contains, a statement about recall of the kind a declaration can legitimately cover.

\paragraph{A filename heuristic.} One attribution of a blind spot to the component it hides is a filename heuristic and is labelled as one: where a route prefix the analysis could not follow has literal segments, a component whose path carries all of them as whole directory names is reported as hidden behind that prefix. It is wrong in two directions that the harness states: components behind a hole whose directory is named differently are reported only as having no node, and a prefix made entirely of parameters is attributed to no component. It declines every one of the \RoundTwoUnexplained{} unaccounted rows in the corpus on its own first two conditions, so it reaches none of the recorded reasons in neither form, and what is recorded here is a property of the rule and not of any figure in Chapter~\ref{ch:results}.

\section{The intermediate must tier, construct by construct}
\label{app:mustledger}

Table~\ref{tab:mustledger} lists every construct that the \MustCtxEdges{} must edges of the intermediate version cite, with the edge that cites it and the condition the current extractor reads at it (Section~\ref{sec:res-round2}). The condition is read by the extractor's own readers, so the table records the extractor's decisions, and a reader checks it by opening each construct at the line given.

\TabMustLedger{The constructs behind the intermediate must tier and the condition the extractor reads
at each. \emph{none} means the construct runs whenever its template or handler does.}
